% Source: Solutions/Pula Karan-PS3-Q6-10.tex (typed LaTeX), Problem 6
% Student's \subsection* headings -> \paragraph.
\submittedby{Pula Karan}

We analyze the subspace topology on $L$ under both product topologies.

\paragraph{Part 1: Subspace Topology inherited from $\mathbb{R}_l\times\mathbb{R}$.}
A basis for $\mathbb{R}_l\times\mathbb{R}$ consists of rectangles of the form
$[a,b)\times(c,d)$ with $a<b$ and $c<d$.

\begin{enumerate}[label=(\alph*)]
  \item \textbf{Vertical Line ($L: x=k$):}
    For any $(k,y_0)\in L$, a basic open neighborhood in $\mathbb{R}_l\times\mathbb{R}$ is
    $U=[k,k+\epsilon)\times(y_0-\epsilon,y_0+\epsilon)$. Intersecting $U$ with $L$ yields:
    \[ U\cap L=\{k\}\times(y_0-\epsilon,y_0+\epsilon) \]
    This traces open intervals along the $y$-axis. Thus, $L$ inherits the
    \textbf{standard Euclidean topology} ($\mathbb{R}$).

  \item \textbf{Horizontal Line ($L: y=k$):}
    For any $(x_0,k)\in L$, a basic open neighborhood in $\mathbb{R}_l\times\mathbb{R}$ is
    $U=[x_0,x_0+\epsilon)\times(k-\epsilon,k+\epsilon)$. Intersecting $U$ with $L$ yields:
    \[ U\cap L=[x_0,x_0+\epsilon)\times\{k\} \]
    This traces right-half-open intervals along the $x$-axis. Thus, $L$ inherits the
    \textbf{lower limit topology} ($\mathbb{R}_l$).

  \item \textbf{Non-Vertical Line ($L: y=mx+k$, $m\neq0$):}
    For any $(x_0,y_0)\in L$, consider basic open sets
    $U=[x_0,x_0+\epsilon)\times(y_0-\delta,y_0+\delta)$.
    \begin{itemize}
      \item If $m>0$, as $x$ increases from $x_0$, $y$ increases from $y_0$. For sufficiently
            small $\delta>0$, $U\cap L$ isolates points corresponding to
            $x\in[x_0,x_0+\delta/m)$.
      \item If $m<0$, as $x$ increases from $x_0$, $y$ decreases from $y_0$. Intersecting
            $[x_0,x_0+\epsilon)$ with $(y_0-\delta,y_0+\delta)$ again isolates points
            corresponding to $x\in[x_0,x_0+\delta/|m|)$.
    \end{itemize}
    In both cases, basic open sets induce half-open intervals $[x_0,x_0+r)$ along $L$. Thus,
    $L$ inherits the \textbf{lower limit topology} ($\mathbb{R}_l$).
\end{enumerate}

\paragraph{Part 2: Subspace Topology inherited from $\mathbb{R}_l\times\mathbb{R}_l$ (Sorgenfrey Plane).}
A basis for $\mathbb{R}_l\times\mathbb{R}_l$ consists of half-open rectangles of the form
$[a,b)\times[c,d)$.

\begin{enumerate}[label=(\alph*)]
  \item \textbf{Vertical ($x=k$), Horizontal ($y=k$), and Positive Slope ($y=mx+k$, $m>0$):}
    For any $(x_0,y_0)\in L$, intersecting $[x_0,x_0+\epsilon)\times[y_0,y_0+\epsilon')$ with
    $L$ isolates points corresponding to $[x_0,x_0+\delta)$ where
    $\delta=\min(\epsilon,\epsilon'/m)$ for $m>0$ (or $\epsilon$ if vertical/horizontal). Thus,
    $L$ inherits the \textbf{lower limit topology} ($\mathbb{R}_l$).

  \item \textbf{Oblique Line with Negative Slope ($y=mx+k$, $m<0$):}
    Let $(x_0,y_0)\in L$ and consider $U=[x_0,x_0+\epsilon)\times[y_0,y_0+\epsilon)$.
    For any $x>x_0$, $y=mx+k<mx_0+k=y_0$, so $(x,y)\notin U$.
    For $x<x_0$, $x\notin[x_0,x_0+\epsilon)$, so $(x,y)\notin U$.
    Therefore, $U\cap L=\{(x_0,y_0)\}$. Since every singleton is open, $L$ inherits the
    \textbf{discrete topology}.
\end{enumerate}
